\section{Conclusion}

Overall, the results are satisfactory. I managed to prototype a pipeline that retrieves YouTube videos, transcribes them in text with an encoder-decoder transformer, and saves the result in a csv file. I retrieved 72 speeches of Giorgia Meloni, for a total of $120,116$ tokens.

I compared them with the De Gasperi corpus, after extracting 473 public speeches for a total of $609,142$ tokens. I calculated the MTLD for both corpora, showing that De Gasperi's speeches are, on average, 30\% more diversified in their lexicon.

Then, I compared the two Prime Ministers' 1000 most-used words: they have 487 words in common, with a Jaccard similarity index of 32\%. The highest-frequency words in common show transversal themes, such as freedom, peace and nation, while the highest-frequency words not in common reflect the PM's different political and societal landscapes. De Gasperi often talked about communism and socialism, while Meloni focuses on Ukraine, Trump and the Euro.

Lastly, I performed a sentiment analysis on each speech using a BERT classifier.  The results underscore the importance of domain-appropriate training data: the model was trained on product reviews, not on political rhetoric, which translated to an average confidence of roughly 40\% per classification.